Autonomous agents built on large language models face a fundamental bottleneck: as tasks extend over long horizons, redundant and low-value information accumulates within the finite context window, progressively degrading reasoning quality and decision accuracy.
This problem demands a principled approach to managing what information occupies the context at each reasoning step.
We present GenericAgent (GA), a general-purpose LLM agent system designed around the principle of \textbf{context information density maximization}, supporting three complementary operating modes: Interactive, Task, and Reflect, for sustained autonomous operation.
GA makes four main contributions:
(1)~an L0--L3 hierarchical memory architecture that loads information on demand and validates updates through action-level verification;
(2)~a reflection-driven self-evolution mechanism that progressively compresses experiential knowledge from natural-language experience to structured SOPs to executable code and reusable skills;
(3)~a minimal set of seven atomic tools whose principled composition achieves full environmental control;
and (4)~a simplified HTML browser control scheme with incremental DOM monitoring for efficient web interaction.
We release GenericAgent as an open-source system to facilitate research on long-term autonomous agents.
